\begin{figure}[H]
\centering
\begin{tikzpicture}[>=Latex,x=1.05cm,y=0.82cm]
    \def\betaval{-0.8}
    \pgfmathsetmacro{\gam}{1/sqrt(1-\betaval*\betaval)}

    \fill[gray!4] (-4.30,-2.15) rectangle (4.30,3.78);

    \begin{scope}
        \clip (-4.30,-2.15) rectangle (4.30,3.78);

        \foreach \x in {-4,-3,...,4}
            \draw[mink grid x] (\x,-2.15) -- (\x,3.78);
        \foreach \y in {-2,-1,...,3}
            \draw[mink grid t] (-4.30,\y) -- (4.30,\y);

        \foreach \i in {-5,-4,...,5}{
            \pgfmathsetmacro{\xp}{0.70*\i}
            \draw[minkdarkred!45,line width=0.34pt]
                ({\gam*(\xp+\betaval*(-6))},{\gam*((-6)+\betaval*\xp)}) --
                ({\gam*(\xp+\betaval*(6))},{\gam*((6)+\betaval*\xp)});
        }
        \foreach \i in {-5,-4,...,5}{
            \pgfmathsetmacro{\ctp}{0.70*\i}
            \draw[minkdarkred!45,line width=0.34pt]
                ({\gam*((-6)+\betaval*\ctp)},{\gam*(\ctp+\betaval*(-6))}) --
                ({\gam*((6)+\betaval*\ctp)},{\gam*(\ctp+\betaval*(6))});
        }

        \draw[minkorange!70,mink dashed] (-3.75,3.75) -- (2.15,-2.15);
        \draw[minkorange!70,mink dashed] (3.75,3.75) -- (-2.15,-2.15);
    \end{scope}

    \draw[mink axis] (-4.30,0) -- (4.40,0)
        node[below] {\contour{white}{$x$}};
    \draw[mink axis] (0,-2.15) -- (0,3.90)
        node[right] {\contour{white}{$ct$}};

    \draw[mink axis prime] (-4.30,3.44) -- (2.68,-2.14)
        node[below right] {\contour{white}{$x'$}};
    \draw[mink axis prime] (1.72,-2.15) -- (-3.03,3.79)
        node[left] {\contour{white}{$ct'$}};

    \coordinate (P) at (1.50,2.25);
    \fill[minkpurple] (P) circle (3.5pt);
    \node[minkpurple,font=\small,above right=2pt] at (P)
        {\contour{white}{$P$}};
    \draw[mink dashed,minkdarkblue] (P) -- (1.50,0);
    \draw[mink dashed,minkdarkblue] (P) -- (0,2.25);

    \node[minkdarkblue,font=\small,align=left,anchor=east] at (4.08,3.47)
        {\contour{white}{$x=100\ \mathrm{km}$}\\
         \contour{white}{$t=0.5\ \mathrm{ms}$}\\
         \contour{white}{$y=10\ \mathrm{km},\ z=1\ \mathrm{km}$}};
    \node[minkdarkred,font=\small,align=left,anchor=west] at (-4.08,-1.55)
        {\contour{white}{$S'$ se mueve con}\\
         \contour{white}{$\beta=-0.8$}};

\end{tikzpicture}
\caption{Diagrama de Minkowski del evento $P$ para una velocidad relativa negativa. La malla azul representa coordenadas constantes de $S$ y la roja las de $S'$; esta última aparece oblicua debido a la transformación de Lorentz.}
\label{fig:taller2-transformacion-1}
\end{figure}